% !TEX TS-program = pdflatex
% !TeX program = pdflatex
% !TEX encoding = UTF-8
% !TEX spellcheck = en_US

\documentclass[xcolor=table]{beamer}

\usepackage{../extra/beamer/karimnlp}
\input{options}

\subtitle[10- Discourse coherence]{Chapter 10\\Discourse coherence} 

\changegraphpath{../img/coherence/}

\begin{document}
	
\begin{frame}
	\frametitle{\inserttitle}
	\framesubtitle{\insertshortsubtitle: Introduction}

	\begin{exampleblock}{Example of a paragraph}
		\begin{center}
			\Large\bfseries
			Winter is one of the four seasons of the year.
			The report was carefully prepared.
			Rain is a common feature of this season.
			I read a report about this season.
		\end{center}
	\end{exampleblock}
	
	\begin{itemize}
		\item What is the main topic of this paragraph?
		\item Is there a coherence problem? If so, describe it.
		\item How could this text be improved for clarity?
	\end{itemize}
	

\end{frame}

\begin{frame}
	\frametitle{\inserttitle}
	\framesubtitle{\insertshortsubtitle}
	
	\begin{itemize}
		\item How can discourse coherence be detected?
		\item \optword{Relations-based coherence}: discourse sentences are bound together by relations.
		\begin{itemize}
			\item Rhetorical Structure Theory (RST) 
			\item Penn Discourse Treebank (PDTB)
		\end{itemize}
		\item \optword{Entity-based coherence}: discourse sentences focus on the same entity or topic.
		\begin{itemize}
			\item Centering Theory 
		\end{itemize}
		\item \optword{Lexical cohesion}: adjacent or related sentences share words or semantic links.
		\begin{itemize}
			\item Lexical similarity measures
		\end{itemize}
	\end{itemize}
	
\end{frame}


\begin{frame}
	\frametitle{\inserttitle}
	\framesubtitle{\insertshortsubtitle: Some humor}

	\begin{center}
		\vgraphpage{humor/humor1.jpg}
	\end{center}

\end{frame}

\begin{frame}
	\frametitle{\inserttitle}
	\framesubtitle{\insertshortsubtitle: Plan}

	\begin{multicols}{2}
	%	\small
	\tableofcontents
	\end{multicols}

\end{frame}

%===================================================================================
\section{Coherence relations}
%===================================================================================

\begin{frame}
	\frametitle{\insertshortsubtitle}
	\framesubtitle{\insertsection}
	
	\begin{itemize}
		\item \textbf{Coherence relations} describe the logical and semantic connections between discourse units (e.g., sentences or clauses).
		
		\item They explain how different parts of a text are linked to form a coherent whole.
		
		\item Common types of coherence relations include:
		\begin{itemize}
			\item \textbf{Cause / Explanation}
			\item \textbf{Contrast}
			\item \textbf{Elaboration}
			\item \textbf{Temporal (sequence of events)}
		\end{itemize}
		
		\item These relations can be:
		\begin{itemize}
			\item \textbf{Explicit}: signaled by discourse markers (e.g., \textit{because, however, then})
			\item \textbf{Implicit}: inferred from context
		\end{itemize}
		
		\item Formal frameworks such as \textbf{RST} and \textbf{PDTB} model these relations.
	\end{itemize}
	
\end{frame}

\subsection{Rhetorical Structure Theory (RST)}

\begin{frame}
	\frametitle{\insertshortsubtitle: \insertsection}
	\framesubtitle{\insertsubsection}
	
	\begin{itemize}
		\item Structure  
		\begin{itemize}
			\item \optword{Nucleus (N)}: a central unit that can be interpreted independently.
			\item \optword{Satellite (S)}: a dependent unit that is interpreted with reference to the nucleus.
		\end{itemize}
		
		\item Communicative perspective
		\begin{itemize}
			\item \optword{Writer}: organizes the text to achieve a communicative goal.
			\item \optword{Reader}: interprets the relations between discourse units.
		\end{itemize}
		
		\item Relation definition
		\begin{itemize}
			\item Constraints on the nucleus
			\item Constraints on the satellite 
			\item Constraints on their combination 
			\item Intended effect on the reader
		\end{itemize}
	\end{itemize}
	
\end{frame}

\begin{frame}
	\frametitle{\insertshortsubtitle: \insertsection}
	\framesubtitle{\insertsubsection: RST relations (Elaboration)}
	
	\begin{itemize}
		\item \optword{Elaboration}: S provides additional information about the situation presented in N.
		\begin{itemize}
			\item \textbf{Constraints on S}: S presents additional information about the situation or some aspect of N.
			\item \textbf{Constraints on N}: N is interpretable independently.
			\item \textbf{Constraints on the combination}: S elaborates on N.
			\item \textbf{Effect}: The reader recognizes S as providing additional information about N.
			\item \textbf{Locus of the effect}: N
			\item E.g. \expword{[\textsubscript{N} The exam is easy.] [\textsubscript{S} It consists of only three short questions.]}
		\end{itemize}
	\end{itemize}
	
\end{frame}

\begin{frame}
	\frametitle{\insertshortsubtitle: \insertsection}
	\framesubtitle{\insertsubsection: RST relations (Evidence)}
	
	\begin{itemize}
		\item \optword{Evidence}: S provides information intended to increase the reader's belief in N.
		\begin{itemize}
			\item \textbf{Constraints on N}: The reader may not believe N to a degree satisfactory to the writer.
			\item \textbf{Constraints on S}: The reader believes S or finds it credible.
			\item \textbf{Constraints on the combination}: The reader's comprehension of S increases belief in N.
			\item \textbf{Effect}: The reader's belief in N is increased.
			\item \textbf{Locus of the effect}: N
			\item E.g. \expword{[\textsubscript{N} Mustafa must be here.] [\textsubscript{S} His car is parked outside.]}
		\end{itemize}
	\end{itemize}
	
\end{frame}

\begin{frame}
	\frametitle{\insertshortsubtitle: \insertsection}
	\framesubtitle{\insertsubsection: RST relations (Contrast)}
	
	\begin{itemize}
		\item \optword{Contrast}: A contrastive relation between two or more nuclei.
		\begin{itemize}
			\item \textbf{Constraints on N}: Each unit is a nucleus (multi-nuclear relation).
			\item \textbf{Constraints on the combination}: The situations presented in the nuclei are 
			(a) similar in several respects, 
			(b) different in some respects, and 
			(c) compared with respect to these differences.
			\item \textbf{Effect}: The reader recognizes both the similarities and the differences between the nuclei.
			\item \textbf{Locus of the effect}: the nuclei
			\item E.g. \expword{[\textsubscript{N} He voted "No" to the new constitution.] [\textsubscript{N} His brother voted "Yes".]}
		\end{itemize}
	\end{itemize}
	
\end{frame}

\begin{frame}
	\frametitle{\insertshortsubtitle: \insertsection}
	\framesubtitle{\insertsubsection: RST relations (Antithesis)}
	
	\begin{itemize}
		\item \optword{Antithesis}: A contrastive relation between N and S in which S is presented so as to increase the reader's positive regard for N.
		\begin{itemize}
			\item \textbf{Constraints on N}: The writer has a positive evaluation of the situation presented in N.
			\item \textbf{Constraints on S}: S presents a situation that contrasts with N.
			\item \textbf{Constraints on the combination}: The situations in N and S are in contrast (cf. Contrast).
			\item \textbf{Effect}: The reader's positive regard for N is increased.
			\item \textbf{Locus of the effect}: N
			\item E.g. \expword{[\textsubscript{N} I defended the republic when I was young;] [\textsubscript{S} I won't give it up now that I'm old.]}
		\end{itemize}
	\end{itemize}
	
\end{frame}


\begin{frame}
	\frametitle{\insertshortsubtitle: \insertsection}
	\framesubtitle{\insertsubsection: Subject matter vs. presentational relations}
	
	\begin{center}
		\footnotesize
	\begin{tblr}{
			colspec = {p{.4\textwidth}lp{.4\textwidth}},
			row{odd} = {lightblue},
			row{even} = {lightyellow},
			row{1} = {darkblue},
			rowsep = 0pt,
		}

		\bfseries\textcolor{white}{Subject matter} && \bfseries\textcolor{white}{Presentational}\\
		
		Elaboration
		
		Circumstance
		
		Solutionhood
		
		Volitional Cause
		
		Volitional Result
		
		Non-Volitional Cause
		
		Non-Volitional Result
		
		Purpose
		
		Condition
		
		Otherwise
		
		Interpretation
		
		Evaluation
		
		Restatement
		
		Summary
		
		Sequence
		
		Contrast
		
		&&
		Motivation
		
		Antithesis
		
		Background
		
		Enablement
		
		Evidence
		
		Justification
		
		Concession\\
	\end{tblr}
	\end{center}
	
\end{frame}

\begin{frame}
	\frametitle{\insertshortsubtitle: \insertsection}
	\framesubtitle{\insertsubsection: semantic vs. pragmatic vs. textual relations}
	
	
	\begin{center}
		\footnotesize
	\begin{tblr}{
			colspec = {p{.3\textwidth}lp{.25\textwidth}lp{.2\textwidth}},
			row{odd} = {lightblue},
			row{even} = {lightyellow},
			row{1} = {darkblue},
			rowsep = 0pt,
		}

		\bfseries\textcolor{white}{Semantic} && \bfseries\textcolor{white}{Pragmatic} && \bfseries\textcolor{white}{Textual} \\
		
		Elaboration
		
		Circumstance
		
		Solutionhood
		
		Volitional Cause
		
		Volitional Result
		
		Non-Volitional Cause
		
		Non-Volitional Result
		
		Purpose
		
		Condition
		
		Interpretation
		
		Evaluation
		
		Sequence
		
		Contrast
		
		&&
		
		Motivation
		
		Antithesis
		
		Enablement
		
		Evidence 
		
		Justification
		
		Concession
		
		&&
		
		Background
		
		Restatement
		
		Summary\\
	\end{tblr}
	\end{center}
	
\end{frame}

\subsection{Penn Discourse Treebank (PDTB)}

\begin{frame}
	\frametitle{\insertshortsubtitle: \insertsection}
	\framesubtitle{\insertsubsection}
	
	\begin{itemize}
		\item Annotation of discourse relations via connectives
		\begin{itemize}
			\item \optword{Explicit connectives}: lexical items that signal a discourse relation
			\begin{itemize}
				\item Subordinating conjunctions (e.g., \expword{'when', 'because', 'although', 'if'})
				\item Coordinating conjunctions (e.g., \expword{'and', 'but', 'or'})
				\item Discourse adverbials (e.g., \expword{'however', 'therefore', 'then'})
				\item Prepositional phrases (e.g., \expword{'as a result', 'in addition', 'in fact'})
			\end{itemize}
			
			\item \optword{Implicit relations}: inferred between adjacent sentences without an explicit connective
		\end{itemize}
		
		\item Arguments annotation
		\begin{itemize}
			\item ARG2: the argument syntactically associated with the connective
			\item ARG1: the other argument (often preceding ARG2)
		\end{itemize}
	\end{itemize}
	
\end{frame}


\begin{frame}
	\frametitle{\insertshortsubtitle: \insertsection}
	\framesubtitle{\insertsubsection: Connectives senses}
	
	\begin{table}
	{\tiny\bfseries
	\begin{tblr}{
			colspec = {p{.1\textwidth}lp{.8\textwidth}},
			row{odd} = {lightblue},
			row{even} = {lightyellow},
			row{1} = {darkblue},
			rowsep = 1pt,
		}

		\bfseries\textcolor{white}{Connective} && \bfseries\textcolor{white}{Senses}\\
		
		after && succession (523), succession-reason (50), other (4) \\
		since && reason (94), succession (78), succession-reason (10), other (2) \\
		when && synchrony (477), succession (157), general (100), succession-reason (65), synchrony-general (50),
		synchrony-reason (39), hypothetical (11), implicit assertion (11), synchrony-hypothetical (10), other
		(69) \\
		while && juxtaposition (182), synchrony (154), contrast (120), expectation (79), opposition (78), conjunction
		(39), synchrony-juxtaposition (26), synchrony-conjunction (21), synchrony-contrast(22), comparison (18), synchrony-opposition (11), other (31) \\
		meanwhile && synchrony-conjunction (92), synchrony (26), conjunction (25), synchrony-juxtaposition (15),
		other(35)\\
		but && Contrast (1609), juxtaposition (636), contra-expectation (494), comparison (260), opposition
		(174), conjunction (63), conjunction-pragmatic contrast (14), pragmatic-contrast (14), other (32)
		however contrast (254), juxtaposition (89), contra-expectation (70), comparison (49), opposition (31),
		other (12)\\
		although && expectation (132), contrast (114) juxtaposition (34), contra-expectation (21), comparison (16),
		opposition (9), other (2)\\
		and && conjunction (2543), list (210), result-conjunction (138), result (38), precedence-conjunction (30),
		juxtaposition (11), other(30)\\
		if && hypothetical (682), general (175), unreal present (122), factual present (73), unreal past (53), expectation (34), implicit assertion (29), relevance (20), other (31)\\
	\end{tblr}}
	\caption{Distribution of senses for selected connectives \cite{2008-prasad-al}}
	\end{table}
	
\end{frame}

\begin{frame}
	\frametitle{\insertshortsubtitle: \insertsection}
	\framesubtitle{\insertsubsection: Connective categories}
	
	\begin{figure}
		\centering\vspace{-12pt}
		\begin{tcolorbox}[colback=white, colframe=blue, boxrule=1pt, text width=.9\textwidth]
		\begin{minipage}{.3\textwidth}
			\tiny\bfseries\vspace{-6pt}
			\begin{itemize}
				\item CONTINGENCY
				\begin{itemize}\tiny\bfseries
					\item Cause
					\begin{itemize}\tiny\bfseries
						\item Reason
						\item Result
					\end{itemize}
					\item Pragmatic Cause
					\begin{itemize}\tiny\bfseries
						\item Justification
					\end{itemize}
					\item Condition
					\begin{itemize}\tiny\bfseries
						\item Hypothetical
						\item General
						\item Unreal Present
						\item Unreal Past
						\item Factual Present
						\item Factual Past
					\end{itemize}
					\item Pragmatic Condition
					\begin{itemize}\tiny\bfseries
						\item Relevance
						\item Implicit Assertion
					\end{itemize}
				\end{itemize}
			\end{itemize}
		\end{minipage}
		\begin{minipage}{.3\textwidth}
			\tiny\bfseries
			\begin{itemize}\tiny\bfseries
				\item TEMPORAL
				\begin{itemize}\tiny\bfseries
					\item Asynchronous
					\item Synchronous
					\begin{itemize}\tiny\bfseries
						\item Precedence
						\item Succession
					\end{itemize}
				\end{itemize}
				\item COMPARISON
				\begin{itemize}\tiny\bfseries
					\item Contrast
					\begin{itemize}\tiny\bfseries
						\item Juxtaposition 
						\item Opposition
					\end{itemize}
					\item Pragmatic Contrast
					\item Concession
					\begin{itemize}\tiny\bfseries
						\item Expectation
						\item Contra-expectation
					\end{itemize}
					\item Pragmatic Concession
				\end{itemize}
			\end{itemize}
		\end{minipage}
		\begin{minipage}{.3\textwidth}
			\tiny\bfseries
			\begin{itemize}\tiny\bfseries
				\item EXPANSION
				\begin{itemize}\tiny\bfseries
					\item Conjunction
					\item Instantiation
					\item Restatement
					\begin{itemize}\tiny\bfseries
						\item Specification
						\item Equivalence
						\item Generalization
					\end{itemize}
					\item Alternative
					\begin{itemize}\tiny\bfseries
						\item Conjunction
						\item Disjunction
						\item Chosen Alternative
					\end{itemize}
					\item Exception
					\item List
				\end{itemize}
			\end{itemize}
		\end{minipage}\vspace{-6pt}
		\end{tcolorbox}\vspace{-6pt}
		\caption{PDTB sense hierarchy \cite{2008-prasad-al}}
	\end{figure}
	
\end{frame}

%===================================================================================
\section{Discourse structure parsing}
%===================================================================================

\begin{frame}
	\frametitle{\insertshortsubtitle}
	\framesubtitle{\insertsection}
	
	\begin{itemize}
		\item A discourse is considered coherent if semantic or rhetorical relations hold between its constituent units.
		\item Discourse structure can be represented using different frameworks:
		\begin{itemize}
			\item \textbf{RST}: relations are organized hierarchically in a tree
			\item \textbf{PDTB}: relations are annotated between pairs of clauses or sentences
		\end{itemize}
	\end{itemize}
	
\end{frame}

\subsection{RST parsing}

\begin{frame}
	\frametitle{\insertshortsubtitle: \insertsection}
	\framesubtitle{\insertsubsection}
	
	\begin{itemize}
		\item Construct a discourse relations tree using RST.
		\item RST parsing typically involves two main steps:
		\begin{enumerate}
			\item Segmentation into Elementary Discourse Units (EDUs);
			\item Classification of relations (including nuclearity and relation type).
		\end{enumerate}
	\end{itemize}
	
	\begin{center}
		\hgraphpage[0.7\textwidth]{RST-tree.pdf}
	\end{center}
	
\end{frame}

\begin{frame}
	\frametitle{\insertshortsubtitle: \insertsection}
	\framesubtitle{\insertsubsection: Elementary Discourse Units (EDUs) detection}
	
	\begin{itemize}
		\item Elementary Discourse Units (EDUs) are the minimal building blocks of discourse.
		\item A single sentence may contain multiple EDUs.
		\item Example: 
		\expword{[Mr. Rambo says]\textsubscript{e1} [that a 3.2-acre property]\textsubscript{e2} [overlooking the San Fernando Valley]\textsubscript{e3} [is priced at \$4 million]\textsubscript{e4} [because the late actor Erroll Flynn once lived there.]\textsubscript{e5}}
		\item EDU segmentation methods:
		\begin{itemize}
			\item \optword{Rule-based}
			\begin{itemize}
				\item Parsing-based: using syntactic parse trees
				\item Surface features: punctuation, lexical cues (e.g., discourse connectives)
			\end{itemize}
			\item \optword{Statistical / Machine Learning}
			\begin{itemize}
				\item \optword{Feature-based}: using syntactic and surface features to predict EDU boundaries
				\item \optword{Embeddings-based}: sequence labeling using neural embeddings. E.g., \href{https://github.com/PKU-TANGENT/NeuralEDUSeg}{NeuralEDUSeg}
			\end{itemize}
		\end{itemize}
	\end{itemize}
	
\end{frame}

\begin{frame}
	\frametitle{\insertshortsubtitle: \insertsection}
	\framesubtitle{\insertsubsection: EDU detection by embeddings}
	
	\begin{figure}
		\centering
		\hgraphpage[.45\textwidth]{EDU_seg.pdf}
		\caption{EDU segmentation as proposed by \cite{2018-wang-al}; figure adapted from \cite{2019-jurafsky-martin}}
	\end{figure}
	
\end{frame}

\begin{frame}
	\frametitle{\insertshortsubtitle: \insertsection}
	\framesubtitle{\insertsubsection: Relations classification}
	
	\begin{minipage}{.65\textwidth}
		\begin{itemize}
			\item The most commonly used method: SHIFT-REDUCE parsing.
			\item Configuration: $C = (\sigma, \beta, A)$
			\begin{itemize}
				\item $\sigma$ is the stack
				\item $\beta$ is the input buffer
				\item $A$ is the list of constructed arcs (relations)
				\item $C_{initial} = (\varnothing, w, \emptyset)$;
				$C_{final} = (\varnothing, \varnothing, A)$
			\end{itemize}
		\end{itemize}
	\end{minipage}
	\begin{minipage}{.33\textwidth}
		\hgraphpage{RST-transitions.pdf}
	\end{minipage}
	
	\begin{itemize}
		\item Operations:
		\begin{itemize}
			\item \optword{Shift}: move the first element of the buffer onto the stack
			\item \optword{Reduce(l, d)}: merge the top two subtrees on the stack
			\begin{itemize}
				\item $l$ = label of the coherence relation
				\item $d \in \{\text{NN}, \text{NS}, \text{SN}\}$ = nuclearity direction
			\end{itemize}
			\item \optword{Pop Root}: remove the top subtree from the stack once fully attached
		\end{itemize}
		\item Training approaches:
		\begin{itemize}
			\item Features-based: using syntactic and surface features to predict transitions
			\item Embeddings-based: using neural embeddings for transition prediction
		\end{itemize}
	\end{itemize}
	
\end{frame}


\begin{frame}
	\frametitle{\insertshortsubtitle: \insertsection}
	\framesubtitle{\insertsubsection: Relations classification (example)}
	
	\begin{figure}
		\begin{minipage}{0.25\textwidth}
			\hgraphpage{RST_exp.pdf}
		\end{minipage}
		\begin{minipage}{0.70\textwidth}
			\scriptsize
			\begin{itemize}
				\item e\textsubscript{1}: American Telephone \& Telegraph Co. said it
				\item e\textsubscript{2}: will lay off 75 to 85 technicians here , effective Nov. 1.
				\item e\textsubscript{3}: The workers install , maintain and repair its private branch exchanges,
				\item e\textsubscript{4}: which are large intracompany telephone networks.
			\end{itemize}
		\end{minipage}
		
		\small
		\begin{tabular}{lllll}
			\hline\hline 
			Step & Stack & Buffer & Action & Added Relation \\
			\hline
			01 & $\emptyset$ & $e_1, e_2, e_3, e_4$ & SH & / \\
			02 & $e_1$ & $e_2, e_3, e_4$ & SH & / \\
			03 & $e_1, e_2$ & $e_3, e_4$ & RD(attr, SN) & $\widehat{e_1 \mathbf{e_2}}$ \\
			04 & $e_{1:2}$ & $e_3, e_4$ & SH & / \\
			05 & $e_{1:2}, e_3$ & $e_4$ & SH & / \\
			06 & $e_{1:2}, e_3, e_4$ & $\emptyset$ & RD(elab, NS) & $\widehat{\mathbf{e_3} e_4}$ \\
			07 & $e_{1:2}, e_{3:4}$ & $\emptyset$ & RD(elab, SN) & $\widehat{e_{1:2}\mathbf{e_{3:4}}}$ \\
			08 & $e_{1:4}$ & $\emptyset$ & PR & / \\
			\hline\hline
		\end{tabular}
		\caption{Example of RST parsing using Shif-Reduce \cite{2018-yu-al}}
	\end{figure}
	
\end{frame}


\begin{frame}
	\frametitle{\insertshortsubtitle: \insertsection}
	\framesubtitle{\insertsubsection: Relations classification using embeddings (1)}
	
	\begin{itemize}
		\item Case study: the encoder-decoder approach of \cite{2018-yu-al}
		\item Implementation: \url{https://github.com/yunan4nlp/NNDisParser}
		\item Encoder 
		\begin{itemize}
			\item Encode each word $w_i$ with its POS tag $t_i$ (concatenation):
			\[
			x_i^w = [embedding(w_i); embedding(t_i)]
			\]
			\item Contextual representation of words:
			\[
			\{h_1^w, h_2^w, \ldots, h_m^w \} = biLSTM(\{x_1^w, x_2^w, \ldots, x_m^w \})
			\]
			\item Represent each EDU $\{w_s, w_{s+1}, \ldots, w_t\}$ by mean pooling:
			\[
			x^e = \frac{1}{t-s+1} \sum_{k=s}^{t} h_k^w
			\]
			\item Contextual representation of sequential EDUs:
			\[
			\{h_1^e, h_2^e, \ldots, h_n^e \} = biLSTM(\{x_1^e, x_2^e, \ldots, x_n^e \})
			\]
		\end{itemize}
	\end{itemize}
	
\end{frame}

\begin{frame}
	\frametitle{\insertshortsubtitle: \insertsection}
	\framesubtitle{\insertsubsection: Relations classification using embeddings (2)}
	
	\begin{itemize}
		\item Decoder 
		\begin{itemize}
			\item A feed-forward neural layer is used to predict the parser action $o$:
			\[
			o = W [h_{s0}^{sbt}; h_{s1}^{sbt}; h_{s2}^{sbt}; h_{q0}^{e}]
			\]
			\item $h_{q0}^{e}$: contextual representation of the first EDU in the buffer (from the encoder)
			\item $h_{si}^{sbt}$: contextual representation of the $i$-th subtree in the stack
			\item The representation of a subtree spanning EDUs $s = \{e_i, \ldots, e_j\}$ is computed as the mean of the embeddings of the included EDUs:
			\[
			h_{s}^{sbt} = \frac{1}{j-i+1} \sum_{k=i}^{j} h_k^e
			\]
		\end{itemize}
	\end{itemize}
	
\end{frame}


\subsection{PDTB parsing}

\begin{frame}
	\frametitle{\insertshortsubtitle: \insertsection}
	\framesubtitle{\insertsubsection}
	
	\begin{itemize}
		\item Determine the discourse relation between two text spans.
		\item Four main steps: 
		\begin{enumerate}
			\item Identify discourse connectives
			\begin{itemize}
				\item Disambiguate the sense of each connective based on context and punctuation
			\end{itemize}
			\item Determine the two argument spans (ARG1 and ARG2) for each connective
			\begin{itemize}
				\item Similar to EDU identification in RST parsing
			\end{itemize}
			\item Label the relation between these spans (explicit connectives)
			\begin{itemize}
				\item The connective itself serves as the label
			\end{itemize}
			\item Assign a relation between every adjacent pair of sentences (implicit connectives)
			\begin{itemize}
				\item For example, the [CLS] token from a pre-trained BERT model can be passed through a feed-forward layer to predict the relation
			\end{itemize}
		\end{enumerate}
	\end{itemize}
	
\end{frame}


\begin{frame}
	\frametitle{\insertshortsubtitle: \insertsection}
	\framesubtitle{\insertsubsection: Some humor}
	
	\begin{center}
		\vgraphpage{humor/humor2.jpg}
	\end{center}
	
\end{frame}

%===================================================================================
\section{Entity-based coherence}
%===================================================================================

\begin{frame}
	\frametitle{\insertshortsubtitle}
	\framesubtitle{\insertsection}
	
	\begin{itemize}
		\item A discourse exhibits entity-based coherence if it consistently focuses on a central entity.
		\item Throughout the discourse, this entity should remain the primary focus.
	\end{itemize}
	
	\begin{exampleblock}{Example of entity-based coherence \cite{2019-jurafsky-martin}}
		\small
		\begin{minipage}{.48\textwidth}\small
			\begin{enumerate}
				\item John went to his favorite music store to buy a piano. [John]
				\item He had frequented the store for many years. [John]
				\item He was excited that he could finally buy a piano. [John]
				\item He arrived just as the store was closing for the day. [John]
			\end{enumerate}
%			\footnotesize{(Coherent: main entity consistently referenced)}
		\end{minipage}
		\begin{minipage}{.48\textwidth}\small
			\begin{enumerate}
				\item John went to his favorite music store to buy a piano. [John]
				\item It was a store John had frequented for many years. [The store]
				\item He was excited that he could finally buy a piano. [John]
				\item It was closing just as John arrived. [The store]
			\end{enumerate}
%			\footnotesize{(Less coherent: focus alternates between entities)}
		\end{minipage}
	\end{exampleblock}
	
\end{frame}


\subsection{Centering theory}

\begin{frame}
	\frametitle{\insertshortsubtitle: \insertsection}
	\framesubtitle{\insertsubsection: Discourse center}
	
	\begin{itemize}
		\item $\mathbf{U_n}$: an utterance
		\item $\mathbf{C_b(U_n)}$: \optword{Backward-looking center}
		\begin{itemize}
			\item the entity currently salient in the discourse following $U_{n-1}$
		\end{itemize}
		\item $\mathbf{C_f(U_n)}$: \optword{Forward-looking centers}
		\begin{itemize}
			\item potentially salient entities in the next utterance
			\item candidates for $C_b(U_{n+1})$
		\end{itemize}
		\item $\mathbf{C_p(U_n) \in C_f(U_n)}$: \optword{Preferred center}
		\begin{itemize}
			\item the candidate most likely to be $C_b(U_{n+1})$
			\item ranked based on grammatical role (e.g. \expword{subject $>$ object $>$ others}), linear order (e.g., \expword{in Arabic, first position carries more weight}), etc.
		\end{itemize}
		\item $C_b(U_n)$ is preferably realized by a pronoun in $U_{n+1}$ to maintain coherence
	\end{itemize}
	
\end{frame}

\begin{frame}
	\frametitle{\insertshortsubtitle: \insertsection}
	\framesubtitle{\insertsubsection: Transitions}
	
	\begin{itemize}
		\item Transitions ranked from most coherent to least coherent.
		\item \optword{CONTINUE}: the speaker talks about an entity and will continue talking about it.
		\item \optword{RETAIN}: the speaker talks about an entity but will switch to another soon.
		\item \optword{SHIFT}: the speaker switches to a new central entity.
		\begin{itemize}
			\item \textbf{Smooth-SHIFT}: after switching, the speaker will return to this entity later.
			\item \textbf{Rough-SHIFT}: after switching, the speaker will not return to this entity.
		\end{itemize}
	\end{itemize}
	
	
	\begin{center}
		\tiny\bfseries
		\begin{tblr}{
			colspec = {p{.2\textwidth}p{.2\textwidth}p{.2\textwidth}},
			row{odd} = {lightblue},
			row{even} = {lightyellow},
			row{1} = {darkblue},
			rowsep = 2pt,
			}
						
			& \bfseries\color{white}$\mathbf{C_b(U_n) = C_b(U_{n-1})}$
						
			OR $\mathbf{C_b(U_n) = NULL}$
			& \bfseries\color{white}$\mathbf{C_b(U_n) \ne C_b(U_{n-1})}$\\
						
			$\mathbf{C_b(U_n) = C_p(U_n)}$ &
			CONTINUE & Smooth-SHIFT\\
						
			$\mathbf{C_b(U_n) \ne C_p(U_n)}$ &
			RETAIN & Rough-SHIFT\\
		\end{tblr}
	\end{center}
\end{frame}


\subsection{Entity Grid model}

\begin{frame}
	\frametitle{\insertshortsubtitle: \insertsection}
	\framesubtitle{\insertsubsection}
	
	\begin{itemize}
		\item A document is represented as a matrix:
		\begin{itemize}
			\item Rows: sentences;
			\item Columns: entities;
			\item Values: grammatical role of the entity; Subject (S), Object (O), Other (X), or missing (\_).
		\end{itemize}
		\item A document is coherent if it follows typical patterns of entity transitions:
		\begin{itemize}
			\item These patterns are captured by the frequencies of grammatical transitions, e.g., \expword{SS, SO, SX, S\_, OS, OO, OX, O\_, ...}
			\item Machine learning models can use these features to judge document coherence.
		\end{itemize}
	\end{itemize}
	
	
\end{frame}


\begin{frame}
	\frametitle{\insertshortsubtitle: \insertsection}
	\framesubtitle{\insertsubsection: Example of document's representation}
	
	\begin{figure}
		\centering
		{\tiny\bfseries
		\begin{minipage}{.8\textwidth}
			\begin{enumerate}
				\item\ [The Justice Department]\textsubscript{s} is conducting an [anti-trust trial]\textsubscript{o} against [Microsoft Corp.]\textsubscript{x} with [evidence]\textsubscript{x} that [the company]\textsubscript{s} is increasingly attempting to crush [competitors]\textsubscript{o}.
				\item\ [Microsoft]\textsubscript{o} is accused of trying to forcefully buy into [markets]\textsubscript{x} where [its own products]\textsubscript{s} are not competitive enough to unseat [established brands]\textsubscript{o}.
				\item\ [The case]\textsubscript{s} revolves around [evidence]\textsubscript{o} of [Microsoft]\textsubscript{s} aggressively pressuring [Netscape]\textsubscript{o} into merging [browser software]\textsubscript{o}.
				\item\ [Microsoft]\textsubscript{s} claims [its tactics]\textsubscript{s} are commonplace and good economically.
				\item\ [The government]\textsubscript{s} may file [a civil suit]\textsubscript{o} ruling that [conspiracy]\textsubscript{s} to curb [competition]\textsubscript{o} through [collision]\textsubscript{x} is [a violation of the Sherman Act]\textsubscript{o}.
				\item\ [Microsoft]\textsubscript{s} continues to show [inscreased earnings]\textsubscript{o} despite [the trial]\textsubscript{x}.
			\end{enumerate}
		\end{minipage}
		
		\color{blue}
		\begin{tabular}{ccccccccccccccccc}
			& \rotatebox[origin=c]{90}{Department} & \rotatebox[origin=c]{90}{Trial} & \rotatebox[origin=c]{90}{Microsoft} &
			\rotatebox[origin=c]{90}{Evidence} & \rotatebox[origin=c]{90}{Competitors} & \rotatebox[origin=c]{90}{Markets} &
			\rotatebox[origin=c]{90}{Products} & \rotatebox[origin=c]{90}{Brands} & \rotatebox[origin=c]{90}{Case} &
			\rotatebox[origin=c]{90}{Netscape} & \rotatebox[origin=c]{90}{Software} & \rotatebox[origin=c]{90}{Tactics} &
			\rotatebox[origin=c]{90}{Government} & \rotatebox[origin=c]{90}{Suit} & \rotatebox[origin=c]{90}{Earnings} & \\
			
			1. & s & o & s & x & o & - & - & - & - & - & - & - & - & - & - & 1. \\
			2. & - & - & o & - & - & x & s & o & - & - & - & - & - & - & - & 2. \\
			3. & - & - & s & o & - & - & - & - & s & o & o & - & - & - & - & 3. \\
			4. & - & - & s & - & - & - & - & - & - & - & - & s & - & - & - & 4. \\
			5. & - & - & - & - & - & - & - & - & - & - & - & - & s & o & - & 5. \\
			6. & - & x & s & - & - & - & - & - & - & - & - & - & - & - & o & 6. \\
			
		\end{tabular}}
		\caption{Example of a document represented using the Entity Grid model \cite{2008-barzilay-lapata}}
		
	\end{figure}
	
\end{frame}

\begin{frame}
	\frametitle{\insertshortsubtitle: \insertsection}
	\framesubtitle{\insertsubsection: Transitions representation}
	
	\begin{itemize}
		\item A document is represented by the probabilities of grammatical transitions.
		\item The probability of a transition is calculated as the ratio of its occurrences to the total number of transitions in the document.
		\item Example: in the previous entity grid, \expword{P(S\_) = 2/75}.
	\end{itemize}
	\begin{figure}
		{\tiny\bfseries
		\begin{tabular}{lllllllllllllllll}
			\hline
			& s s & s o & s x & s - & o s & o o & o x & o - & x s & x o & x x & x - & - s & - o & - x & - - \\
			\hline
			d1 & .01 & .01 & .00 & .08 & .01 & .00 & .00 & .09 & .00 & .00 & .00 & .03 & .05 & .07 & .03 & .59 \\
			d2 & .02 & .01 & .01 & .02 & .00 & .07 & .00 & .02 & .14 & .14 & .06 & .04 & .03 & .07 & .01 & .36 \\
			d3 & .02 & .00 & .00 & .03 & .09 & .00 & .09 & .06 & .00 & .00 & .00 & .05 & .03 & .07 & .17 & .39 \\
			\hline
		\end{tabular}}
		\caption{Example of a document's representation using grammatical transitions \cite{2008-barzilay-lapata}}
		
	\end{figure}
	
\end{frame}

\begin{frame}
	\frametitle{\insertshortsubtitle: \insertsection}
	\framesubtitle{\insertsubsection: Training}
	
	\begin{itemize}
		\item The algorithm learns to score a text based on its coherence.
		\item Incoherent texts should receive lower scores than coherent ones.
		\item Training data:
		\begin{itemize}
			\item Using texts annotated for coherence by human experts.
			\item Using self-supervised methods by generating incoherent texts from coherent ones:
			\begin{itemize}
				\item \optword{Sentence order discrimination}: compare the original document to a random permutation of its sentences; the model should score the original higher.
				\item \optword{Sentence insertion}: insert a sentence at a random position to create an incoherent version.
				\item \optword{Sentence order reconstruction}: randomly shuffle the sentences of a document, then train the model to restore the correct order.
			\end{itemize}
		\end{itemize}
	\end{itemize}
	
	
\end{frame}

\begin{frame}
	\frametitle{\insertshortsubtitle: \insertsection}
	\framesubtitle{\insertsubsection: Some humor}
	
	\begin{center}
		\vgraphpage{humor/humor3.png}
	\end{center}
	
\end{frame}


\insertbibliography{NLP10}{*}

\end{document}

